\documentclass[addpoints,spanish, 12pt,a4paper]{exam}
%\documentclass[answers, spanish, 12pt,a4paper]{exam}
%\printanswers
\pointpoints{punto}{puntos}
\hpword{Puntos:}
\vpword{Puntos:}
\htword{Total}
\vtword{Total}
\hsword{Resultado:}
\hqword{Ejercicio:}
\vqword{Ejercicio:}

\usepackage[utf8x]{inputenc}
\usepackage[spanish]{babel}
\usepackage{eurosym}
%\usepackage[spanish,es-lcroman, es-tabla, es-noshorthands]{babel}


\usepackage{pgf,tikz}
\usetikzlibrary{shapes, calc, shapes, arrows, math, babel}

\usepackage[margin=1in]{geometry}
\usepackage{amsmath,amssymb}
\usepackage{multicol}
\usepackage{yhmath}
\usepackage{pdflscape}

\pointsinrightmargin % Para poner las puntuaciones a la derecha. Se puede cambiar. Si se comenta, sale a la izquierda.
\extrawidth{-2.4cm} %Un poquito más de margen por si ponemos textos largos.
\marginpointname{ \emph{\points}}

\usepackage{graphicx}
\graphicspath{{../img/}} 

\newcommand{\class}{1º Bachillerato Ciencias}
\newcommand{\examdate}{\today}
\newcommand{\examnum}{Parcial 3ª evaluación}
\newcommand{\tipo}{A}


\newcommand{\timelimit}{50 minutos}

\renewcommand{\solutiontitle}{\noindent\textbf{Solución:}\enspace}


\pagestyle{head}
\firstpageheader{\includegraphics[width=0.2\columnwidth]{header_left}}{\textbf{Departamento de Matemáticas\linebreak \class}\linebreak \examnum}{\includegraphics[width=0.1\columnwidth]{header_right}}
\runningheader{\class}{\examnum}{Página \thepage\ of \numpages}
\runningheadrule


\begin{document}

\noindent
\begin{tabular*}{\textwidth}{l @{\extracolsep{\fill}} r @{\extracolsep{6pt}} }
\textbf{Nombre:} \makebox[3.5in]{\hrulefill} & \textbf{Fecha:}\makebox[1in]{\hrulefill} \\
 & \\
\textbf{Tiempo: \timelimit} & Tipo: \tipo 
\end{tabular*}
\rule[2ex]{\textwidth}{2pt}
Esta prueba tiene \numquestions\ ejercicios. La puntuación máxima es de \numpoints. 
La nota final de la prueba será la parte proporcional de la puntuación obtenida sobre la puntuación máxima. 

\begin{center}


\addpoints
 %\gradetable[h][questions]
	\pointtable[h][questions]
\end{center}

\noindent
\rule[2ex]{\textwidth}{2pt}

\begin{questions}

\question Una asociación de consumidores ha realizado una prueba sobre la duración de unas bombillas de una conocida marca. Ha mantenido encendidas ininterrumpidamente 100 bombillas  hasta que se han estropeado. Sus resultados han sido:



\begin{tabular}{cc}
\hline
   Duración en días &   Nº de bombillas \\
\hline
     $\left[36, 42\right)$ &     12      \\
     $\left[42, 48\right)$ &     28      \\
     $\left[48, 54\right)$ &     45      \\
     $\left[54, 60\right)$ &     15      \\
\hline
\end{tabular}

\begin{parts}
\part[1] Realizar una tabla de frecuencias con los datos que vayas a necesitar para resolver el ejercicio

\begin{solution}\\

\begin{tabular}{rrrrrrrrrr}
\hline
    &   lim\_inf &   lim\_sup &   x\_i &   f\_i &   F\_i &   h\_i &    H\_i &   x\_if\_i &   x\^{}2\_if\_i \\
\hline
  0 &        36 &        42 &    39 &    12 &    12 &  0.12 &   0.12 &      468 &      18252 \\
  1 &        42 &        48 &    45 &    28 &    40 &  0.28 &   0.4  &     1260 &      56700 \\
  2 &        48 &        54 &    51 &    45 &    85 &  0.45 &   0.85 &     2295 &     117045 \\
  3 &        54 &        60 &    57 &    15 &   100 &  0.15 &   1    &      855 &      48735 \\
\hline
  4 &       nan &       nan &   nan &   100 &   nan &  1    & nan    &     4878 &     240732 \\
\hline
\end{tabular}

\end{solution}
\part[1] Calcula la media y la varianza.
\begin{solution}
{'media': 48.78,
 'varianza': 27.83159999999998,
 'desviación típica': 5.27556632031102}
\end{solution}
\part[1] Indica razonadamente en qué intervalos se encuentra la moda y la mediana respectivamente. 
\begin{solution}
\end{solution}
\part[1] Calcula la mediana. Ayuda:
$$P_k=L_i + \frac{k\frac{N}{100}-F_{i-1}}{f_i}\cdot C_i$$
\begin{solution}
 ${'k': 50, 'N': 100.0, 'L_i': 48.0, 'f_i': 45.0, 'F_{i-1}': 40.0, 'C_i': 6.0}$ \\
 49.333333333333336
\end{solution}
\part[1] En el anuncio de televisión el fabricante asegura que sus bombillas duran más de 1200 horas. ¿Qué porcentaje de las bombillas no cumple lo enunciado?
\begin{solution}
${'valor': 50.0, 'N': 100.0, 'L_i': 48.0, 'f_i': 45.0, 'F_{i-1}': 40.0, 'C_i': 6.0}$\\
52.6041666666668
\end{solution}

\end{parts}

\addpoints

        \question En Utebo, el 50\% de sus habitantes mayores de edad tiene más de 60 años, el 40\% de los mayores de edad fuma y el 60\% de los mayores de edad
fuma o es mayor de 60 años. 
Calcula la probabilidad de los siguientes sucesos:
        
        \begin{parts} \part[1] Ser mayor de 60 años y fumar   \begin{solution}  $ \left[  \frac{3}{10}\right] $  \end{solution}
        \part[1] No fumar  \begin{solution}  $ \left[ \frac{3}{5}\right] $  \end{solution}
        \end{parts}
        
        
        \question Dos fábricas producen el mismo televisor. La fábrica A produce el 60\% de todos los televisores.
El 1\% de todos los televisores producidos por la fábrica A salen defectuosos, mientras que el 2\% de los
televisores producidos por la fábrica B son defectuosos. Se selecciona un televisor al azar de entre
todos los fabricados:


        \begin{parts} \part[1] Calcular la probabilidad de que sea defectuoso
  \begin{solution}  $ \left[ \frac{7}{500}, \  \right] $  \end{solution}
\part[1] Si sabemos que el televisor es defectuoso, calcula la probabilidad de que haya sido producido por la
planta A.  \begin{solution}  $ \left[  \frac{3}{7}\right] $  \end{solution}
        \end{parts}

\question Un jugador de baloncesto tiene un porcentaje de acierto en tiros de 2 del 70 \%. Si tira 5
veces:\\ 
    

        \begin{parts} 
        %\part[1] Describe la variable del ejercicio  \begin{solution}   $\left\{ 0 : 0.00243, \  1 : 0.02835, \  2 : 0.1323, \  3 : 0.3087, \  4 : 0.36015, \  5 : 0.16807\right\}$   \end{solution} 
        \part[1] Calcula la probabilidad de que enceste 4  \begin{solution}   $P{\left(X = 4 \right)}=0.36015$   \end{solution} \part[1] Calcula la probabilidad de que enceste al menos 1  \begin{solution}   $P{\left(X \geq 1 \right)}=0.99757$   \end{solution} \part[1] Calcula la probabilidad de que enceste más de 4  \begin{solution}   $P{\left(X > 4 \right)}=0.16807$   \end{solution}
        \end{parts}

\question La duración media de un televisor es de 10 años, con una desviación típica igual a 2 años. Si la
vida útil del electrodoméstico se distribuye normalmente, halla la probabilidad de que al comprar
un televisor: 
    
        
        \begin{parts}  \part[1]  Este dure más de 12 años  \begin{solution}   $P{\left(X \geq 12 \right)}=0.158655253931457$   \end{solution} \part[2]  Este dure entre 8 y 12 años  \begin{solution}   $P{\left(X \geq 8 \wedge X \leq 12 \right)}=0.682689492137086$   \end{solution}
        \end{parts}
        

\addpoints
        
\end{questions}


\newgeometry{left=1 cm,bottom=2cm}
\begin{landscape}
\begin{table}[]
\Large
\centering
\caption{Extracto de tabla de probabilidades de la \textbf{normal estándar $Z(0,1)$}}
\label{my-label}

\begin{tabular}{l|llllllllll}
z   & 0       & 0,01    & 0,02    & 0,03    & 0,04    & 0,05    & 0,06    & 0,07    & 0,08    & 0,09    \\
\hline
0   & 0,5     & 0,50399 & 0,50798 & 0,51197 & 0,51595 & 0,51994 & 0,52392 & 0,5279  & 0,53188 & 0,53586 \\
0,1 & 0,53983 & 0,5438  & 0,54776 & 0,55172 & 0,55567 & 0,55962 & 0,56356 & 0,56749 & 0,57142 & 0,57535 \\
0,2 & 0,57926 & 0,58317 & 0,58706 & 0,59095 & 0,59483 & 0,59871 & 0,60257 & 0,60642 & 0,61026 & 0,61409 \\
0,3 & 0,61791 & 0,62172 & 0,62552 & 0,6293  & 0,63307 & 0,63683 & 0,64058 & 0,64431 & 0,64803 & 0,65173 \\
0,4 & 0,65542 & 0,6591  & 0,66276 & 0,6664  & 0,67003 & 0,67364 & 0,67724 & 0,68082 & 0,68439 & 0,68793 \\
0,5 & 0,69146 & 0,69497 & 0,69847 & 0,70194 & 0,7054  & 0,70884 & 0,71226 & 0,71566 & 0,71904 & 0,7224  \\
0,6 & 0,72575 & 0,72907 & 0,73237 & 0,73565 & 0,73891 & 0,74215 & 0,74537 & 0,74857 & 0,75175 & 0,7549  \\
0,7 & 0,75804 & 0,76115 & 0,76424 & 0,7673  & 0,77035 & 0,77337 & 0,77637 & 0,77935 & 0,7823  & 0,78524 \\
0,8 & 0,78814 & 0,79103 & 0,79389 & 0,79673 & 0,79955 & 0,80234 & 0,80511 & 0,80785 & 0,81057 & 0,81327 \\
0,9 & 0,81594 & 0,81859 & 0,82121 & 0,82381 & 0,82639 & 0,82894 & 0,83147 & 0,83398 & 0,83646 & 0,83891 \\
1   & 0,84134 & 0,84375 & 0,84614 & 0,84849 & 0,85083 & 0,85314 & 0,85543 & 0,85769 & 0,85993 & 0,86214 \\
%1,1 & 0,86433 & 0,8665  & 0,86864 & 0,87076 & 0,87286 & 0,87493 & 0,87698 & 0,879   & 0,881   & 0,88298 \\
%1,2 & 0,88493 & 0,88686 & 0,88877 & 0,89065 & 0,89251 & 0,89435 & 0,89617 & 0,89796 & 0,89973 & 0,90147 \\
%1,3 & 0,9032  & 0,9049  & 0,90658 & 0,90824 & 0,90988 & 0,91149 & 0,91309 & 0,91466 & 0,91621 & 0,91774 \\
%1,4 & 0,91924 & 0,92073 & 0,9222  & 0,92364 & 0,92507 & 0,92647 & 0,92785 & 0,92922 & 0,93056 & 0,93189 \\
%1,5 & 0,93319 & 0,93448 & 0,93574 & 0,93699 & 0,93822 & 0,93943 & 0,94062 & 0,94179 & 0,94295 & 0,94408 \\
%1,6 & 0,9452  & 0,9463  & 0,94738 & 0,94845 & 0,9495  & 0,95053 & 0,95154 & 0,95254 & 0,95352 & 0,95449 \\
%1,7 & 0,95543 & 0,95637 & 0,95728 & 0,95818 & 0,95907 & 0,95994 & 0,9608  & 0,96164 & 0,96246 & 0,96327 \\
%1,8 & 0,96407 & 0,96485 & 0,96562 & 0,96638 & 0,96712 & 0,96784 & 0,96856 & 0,96926 & 0,96995 & 0,97062 \\
%1,9 & 0,97128 & 0,97193 & 0,97257 & 0,9732  & 0,97381 & 0,97441 & 0,975   & 0,97558 & 0,97615 & 0,9767  \\
%2   & 0,97725 & 0,97778 & 0,97831 & 0,97882 & 0,97932 & 0,97982 & 0,9803  & 0,98077 & 0,98124 & 0,98169 \\
%2,1 & 0,98214 & 0,98257 & 0,983   & 0,98341 & 0,98382 & 0,98422 & 0,98461 & 0,985   & 0,98537 & 0,98574 \\
%2,2 & 0,9861  & 0,98645 & 0,98679 & 0,98713 & 0,98745 & 0,98778 & 0,98809 & 0,9884  & 0,9887  & 0,98899 \\
%2,3 & 0,98928 & 0,98956 & 0,98983 & 0,9901  & 0,99036 & 0,99061 & 0,99086 & 0,99111 & 0,99134 & 0,99158 \\
%2,4 & 0,9918  & 0,99202 & 0,99224 & 0,99245 & 0,99266 & 0,99286 & 0,99305 & 0,99324 & 0,99343 & 0,99361 \\
%2,5 & 0,99379 & 0,99396 & 0,99413 & 0,9943  & 0,99446 & 0,99461 & 0,99477 & 0,99492 & 0,99506 & 0,9952  \\
%2,6 & 0,99534 & 0,99547 & 0,9956  & 0,99573 & 0,99585 & 0,99598 & 0,99609 & 0,99621 & 0,99632 & 0,99643 \\
%2,7 & 0,99653 & 0,99664 & 0,99674 & 0,99683 & 0,99693 & 0,99702 & 0,99711 & 0,9972  & 0,99728 & 0,99736 \\
%2,8 & 0,99744 & 0,99752 & 0,9976  & 0,99767 & 0,99774 & 0,99781 & 0,99788 & 0,99795 & 0,99801 & 0,99807 \\
%2,9 & 0,99813 & 0,99819 & 0,99825 & 0,99831 & 0,99836 & 0,99841 & 0,99846 & 0,99851 & 0,99856 & 0,99861 \\
%3   & 0,99865 & 0,99869 & 0,99874 & 0,99878 & 0,99882 & 0,99886 & 0,99889 & 0,99893 & 0,99896 & 0,999   \\
%3,1 & 0,99903 & 0,99906 & 0,9991  & 0,99913 & 0,99916 & 0,99918 & 0,99921 & 0,99924 & 0,99926 & 0,99929 \\
%3,2 & 0,99931 & 0,99934 & 0,99936 & 0,99938 & 0,9994  & 0,99942 & 0,99944 & 0,99946 & 0,99948 & 0,9995  \\
%3,3 & 0,99952 & 0,99953 & 0,99955 & 0,99957 & 0,99958 & 0,9996  & 0,99961 & 0,99962 & 0,99964 & 0,99965 \\
%3,4 & 0,99966 & 0,99968 & 0,99969 & 0,9997  & 0,99971 & 0,99972 & 0,99973 & 0,99974 & 0,99975 & 0,99976 \\
%3,5 & 0,99977 & 0,99978 & 0,99978 & 0,99979 & 0,9998  & 0,99981 & 0,99981 & 0,99982 & 0,99983 & 0,99983 \\
%3,6 & 0,99984 & 0,99985 & 0,99985 & 0,99986 & 0,99986 & 0,99987 & 0,99987 & 0,99988 & 0,99988 & 0,99989 \\
%3,7 & 0,99989 & 0,9999  & 0,9999  & 0,9999  & 0,99991 & 0,99991 & 0,99992 & 0,99992 & 0,99992 & 0,99992 \\
%3,8 & 0,99993 & 0,99993 & 0,99993 & 0,99994 & 0,99994 & 0,99994 & 0,99994 & 0,99995 & 0,99995 & 0,99995 \\
%3,9 & 0,99995 & 0,99995 & 0,99996 & 0,99996 & 0,99996 & 0,99996 & 0,99996 & 0,99996 & 0,99997 & 0,99997 \\
%4   & 0,99997 & 0,99997 & 0,99997 & 0,99997 & 0,99997 & 0,99997 & 0,99998 & 0,99998 & 0,99998 & 0,99998
\end{tabular}
\end{table}
\end{landscape}
\restoregeometry

\end{document}
\grid
